\documentclass[tikz,border=1cm]{standalone}
\usepackage{xcolor}
\usepackage{amsmath,amssymb}

\usepackage[utf8]{inputenc}
\usepackage[spanish]{babel}
\usepackage[
    letterpaper,
    top=1.5cm,
    bottom=2cm,
    left=2cm,
    right=2cm
]{geometry}

\usepackage{graphicx}
\usepackage{fancyhdr}
\usepackage[table,xcdraw]{xcolor}
\usepackage{tikz}
\usetikzlibrary{shapes.geometric, arrows, positioning, fit, backgrounds, calc, babel}
\usepackage{ifpdf}
\usepackage{blindtext}
\usepackage[cmex10]{amsmath}
\usepackage{amssymb}
\usepackage{amsfonts}
\usepackage{float}
\usepackage{siunitx}
\usepackage{listings}
\usepackage{hyperref}
\usepackage{subcaption}
\usepackage{booktabs}
\usepackage{url}
\usepackage{wrapfig}
\usepackage{titlesec}
\usepackage[export]{adjustbox}
\usepackage{imakeidx}
\usepackage{parskip}
\usepackage[numbers]{natbib}

\usepackage{ragged2e}
\usepackage{setspace}
\usepackage{array}
\usepackage{longtable}

% Definición de un nuevo tipo de columna para evitar "Underfull \hbox" en tablas
\newcolumntype{L}[1]{>{\RaggedRight\arraybackslash}p{#1}}

\hypersetup{colorlinks=false,bookmarksopen=true,linkbordercolor={1 1 1}}

\newcommand{\rfigura}[1]{Fig.\ref{#1}}
\newcommand{\rtabla}[1]{Tabla \ref{#1}}

% Datos del informe
\newcommand{\Curso}{IEE2913 --- Diseño Eléctrico (Capstone)}
\newcommand{\TituloInforme}{Informe de Avance 1 (5\%)}
\newcommand{\NumeroGrupo}{10}
\newcommand{\NombreProyecto}{SupaClock: Wearable Biométrico Modular}
\newcommand{\Integrantes}{Tomás Avendaño, Benjamín Sepúlveda, Pablo Uribe}
\newcommand{\FechaEntrega}{7 de abril de 2026}

\lstset{
  basicstyle=\ttfamily\small,
  breaklines=true,
  breakatwhitespace=true,
  postbreak=\mbox{\textcolor{red}{$\hookrightarrow$}\space},
  captionpos=b,
  frame=single,
  numbers=left,
  numberstyle=\tiny\color{gray},
  language=C,
  keywordstyle=\color{blue}\bfseries,
  commentstyle=\color{gray}\itshape,
  stringstyle=\color{red!70!black}
}


\begin{document}
\begin{tikzpicture}[
            node distance=1.5cm and 2.0cm,
            task/.style={rectangle, rounded corners, draw=black, thick, fill=blue!15, text width=3.0cm, align=center, minimum height=1.5cm, font=\small\bfseries},
            queue/.style={rectangle, draw=black!70, dashed, fill=yellow!15, text width=3.5cm, align=center, minimum height=1cm, font=\footnotesize},
            hw/.style={rectangle, draw=black, thick, fill=gray!20, text width=2.5cm, align=center, minimum height=1.0cm, font=\footnotesize},
            arrow/.style={thick, ->, >=stealth},
            data/.style={thick, ->, >=stealth, blue!70, dashed}
            ]
            
            % Tasks
            \node[task] (tsensor) {TaskSensor\\(Prioridad Alta)};
            \node[task, right=3.5cm of tsensor] (tdisplay) {TaskDisplay\\(Prioridad Media)};
            \node[task, below=2cm of tsensor, fill=purple!10, dashed] (tble) {TaskBLE\\(Inactiva/Futuro)};
            
            % Queues
            \node[queue] (qdata) at ($(tsensor)!0.5!(tdisplay)$) {xQueue\\(Datos de Sensores)};
            
            % HW
            \node[hw, above=1.5cm of tsensor] (i2c_hw) {Bus I2C\\(ej. MAX30205)};
            \node[hw, above=1.5cm of tdisplay] (spi_hw) {Bus SPI\\(ST7789)};
            
            % Connections
            \draw[arrow] (i2c_hw) -- node[left, font=\scriptsize] {Polling 50Hz} (tsensor);
            \draw[data] (tsensor) -- node[above, font=\scriptsize] {Temp} (qdata);
            \draw[data] (qdata) -- node[above, font=\scriptsize] {Lectura} (tdisplay);
            \draw[arrow] (tdisplay) -- node[right, font=\scriptsize] {LVGL Flush} (spi_hw);
            
            \draw[data, bend right] (tsensor.south) to node[right, font=\scriptsize] {Datos (Futuro)} (tble.north);
            
        \end{tikzpicture}
\end{document}
